\documentclass[14pt,dvipdfmx]{jsarticle}
\input{lecture-preamble.tex}

% 例題・練習の見出し: \reidai{番号}，\renshu{番号}
\newcommand{\reidai}[1]{\par\noindent\textbf{例題\,#1}\quad}
\newcommand{\renshu}[1]{\par\noindent\textbf{練習\,#1}\quad}

% 指数関数の特徴の図: \expfeature{底}{p}{q}{a>1 なら 1，0<a<1 なら 0}（0<a<1 では a^p，a^q を y 軸の右に置く）
\newcommand{\expfeature}[4]{%
  \begin{tikzpicture}[x=0.45cm, y=0.36cm]
    \draw[-{Stealth}] (-3,0) -- (3,0) node[right]{$x$};
    \draw[-{Stealth}] (0,-0.4) -- (0,4.2) node[above]{$y$};
    \node[below left, font=\footnotesize] at (0,0) {O};
    \ifnum#4=1
      \draw[thick, domain=-3:2.15, samples=60] plot (\x, {#1^(\x)});
    \else
      \draw[thick, domain=-2.15:3, samples=60] plot (\x, {#1^(\x)});
    \fi
    \ifnum#4=1 \def\lblside{left}\else\def\lblside{right}\fi
    \draw[dashed] (#2,0) node[below, font=\footnotesize]{$p$} -- (#2,{#1^(#2)}) -- (0,{#1^(#2)})
      node[\lblside, font=\footnotesize]{$a^p$};
    \draw[dashed] (#3,0) node[below, font=\footnotesize]{$q$} -- (#3,{#1^(#3)}) -- (0,{#1^(#3)})
      node[\lblside, font=\footnotesize]{$a^q$};
    \node[above \lblside, inner sep=1pt, font=\footnotesize] at (0,1) {$1$};
  \end{tikzpicture}%
}

\begin{document}

\lectureheader{数学Ⅱ}{5}{指数関数と対数関数}{1}{指数関数}{2-B}{指数関数の特徴}
  {158}{指数関数の特徴を使って，大小比較について考えよう}

% 同じ節の2枚目以降は，前回までの番号を引き継ぐ
\setcounter{definition}{2}
\setcounter{theorem}{1}
\setcounter{technique}{1}

% ============ 1ページ目（表・左） ============
\heading{〜前時の復習〜}
\begin{theorembox}{指数関数 $y=a^x$ のグラフの特徴}
  \noindent
  \begin{minipage}[t]{0.48\linewidth}
    \textbf{(i)\enspace\boldmath$a>1$ のとき}\par
    \centering
    \begin{tikzpicture}[x=0.5cm, y=0.38cm]
      \draw[-{Stealth}] (-3,0) -- (2.6,0) node[right]{$x$};
      \draw[-{Stealth}] (0,-0.6) -- (0,4.4) node[above]{$y$};
      \node[below left, font=\footnotesize] at (0,0) {O};
      \ifshowanswer
        \begin{scope}[blue]
          \draw[thick, domain=-3:1.75, samples=60] plot (\x, {2.2^(\x)});
          \draw[dashed] (1,0) node[below, font=\footnotesize]{$1$} -- (1,2.2) -- (0,2.2) node[left, font=\footnotesize]{$a$};
          \node[above left, inner sep=1pt, font=\footnotesize] at (0,1) {$1$};
        \end{scope}
      \fi
    \end{tikzpicture}
  \end{minipage}\hfill
  \begin{minipage}[t]{0.48\linewidth}
    \textbf{(ii)\enspace\boldmath$0<a<1$ のとき}\par
    \centering
    \begin{tikzpicture}[x=0.5cm, y=0.38cm]
      \draw[-{Stealth}] (-2.6,0) -- (3,0) node[right]{$x$};
      \draw[-{Stealth}] (0,-0.6) -- (0,4.4) node[above]{$y$};
      \node[below left, font=\footnotesize] at (0,0) {O};
      \ifshowanswer
        \begin{scope}[blue]
          \draw[thick, domain=-1.75:3, samples=60] plot (\x, {0.45^(\x)});
          \draw[dashed] (1,0) node[below, font=\footnotesize]{$1$} -- (1,0.45) -- (0,0.45) node[left, font=\footnotesize]{$a$};
          \node[above right, inner sep=1pt, font=\footnotesize] at (0,1) {$1$};
        \end{scope}
      \fi
    \end{tikzpicture}
  \end{minipage}
  \par\vspace{0.5em}
  \begin{enumerate}[label=\textbf{\arabic*}, leftmargin=1.5zw, itemsep=0.2em]
    \item 2点 $(0,\ \fitblankbf{1})$，$(1,\ \fitblankbf{a})$ を通る．
    \item $\fitblankbf{x \text{軸}}$ を漸近線としてもつ．
    \item $a>1$ のとき\fitblankbf{\text{右上がり}}の曲線，\\
          $0<a<1$ のとき\fitblankbf{\text{右下がり}}の曲線である．
  \end{enumerate}
\end{theorembox}

\vfill
\nextpage

% ============ 2ページ目（表・右） ============
\heading{指数関数の特徴}
\begin{definitionbox}{増加関数・減少関数}
  $x$ の値が増加すると $y$ の値も増加する関数を\fitblankbf{\text{増加関数}}といい，\\
  $x$ の値が増加すると $y$ の値は減少する関数を\fitblankbf{\text{減少関数}}という．
\end{definitionbox}

左のグラフから，$y=a^x$ が増加関数になるか減少関数になるかを考えてみよう．
\par\vspace{0.5em}
\hspace*{1zw}$a>1$ のとき，$y=a^x$ は\fitblankbf{\text{増加関数}}である．\par\vspace{0.6em}
\hspace*{1zw}$0<a<1$ のとき，$y=a^x$ は\fitblankbf{\text{減少関数}}である．
\par\vspace{0.5em}
\begin{theorembox}{指数関数 $y=a^x$ の特徴}
  \begin{enumerate}[label=\textbf{\arabic*}, leftmargin=1.5zw, itemsep=0.3em]
    \item 定義域は実数全体，値域は正の数全体である．
    \item $a>1$ のとき，\fitblankbf{\text{増加関数}}である．\\
          すなわち\quad $p<q \iff a^p \fitblankbf{<} a^q$
    \item $0<a<1$ のとき，\fitblankbf{\text{減少関数}}である．\\
          すなわち\quad $p<q \iff a^p \fitblankbf{>} a^q$
  \end{enumerate}
  \vspace{0.3em}
  \noindent
  \begin{minipage}[t]{0.48\linewidth}
    \textbf{\boldmath$a>1$ のとき}\par\centering
    \expfeature{2}{1}{2}{1}
  \end{minipage}\hfill
  \begin{minipage}[t]{0.48\linewidth}
    \textbf{\boldmath$0<a<1$ のとき}\par\centering
    \expfeature{0.5}{-2}{-1}{0}
  \end{minipage}
\end{theorembox}

\begin{techniquebox}{数の大小比較}
  \begin{itemize}[leftmargin=1.5zw, itemsep=0.2em]
    \item 底を\fitblankbf{\text{そろえて}}，指数の大小を比べる．
    \item 底が \fitblankbf{1} より大きいか小さいかで，大小の向きが決まる．
  \end{itemize}
\end{techniquebox}

\vfill
\nexttime{指数関数を含む方程式・不等式について考えてみよう！}

\end{document}
